\documentclass[11pt]{article}

\usepackage[margin=1in]{geometry}
\usepackage{parskip}
\usepackage{enumitem}
\usepackage{hyperref}

\setlist[itemize]{leftmargin=1.4em, itemsep=2pt, topsep=2pt, parsep=0pt}

\title{\textbf{Summary of Revisions}\\[0.35em]
\large Chaining Tasks, Redefining Work: A Theory of AI Automation}
\author{}
\date{\today}

\begin{document}
\maketitle
\vspace{-2.2em}

\noindent This note lists the main changes in the current draft relative to the previous version, organized by section.

\section*{Abstract}
\begin{itemize}
\item Unchanged from the previous version.
\end{itemize}

\section*{Introduction and Conclusion}
\begin{itemize}
\item Minor revisions only. [Still needs a full pass after we decide on the core content. Peyman will revise both later.]
\item Removed the most direct job design content that the previous iteration carried in these sections.
\item Updated some of the paragraphs so they align with the updated model and discussion sections.
\end{itemize}

\section*{Model}
\begin{itemize}
\item \textbf{Steps and Tasks.} Essentially unchanged. Added a notation summary table and a new illustration of how steps are mapped into tasks (and removed the mapping from tasks to jobs).
\item \textbf{Job design moved out of the main body.} The idea of having jobs is retained: a worker is paid for the duration of their employment, at a rate proportional to the skill the job requires. That, in essence, is what a job is. But \textbf{none} of the surrounding machinery is kept: the endogenous job boundaries, the hand-off and coordination costs, and the optimization of how jobs are designed. I moved all of it to an appendix, to be revisited later, so none of it appears in the main body.
\item \textbf{Firm's cost-minimization problem, restated over two horizons.} With job design no longer in the body, the firm's problem is presented across two horizons:
  \begin{itemize}
  \item \emph{Short run:} with AI available, the firm optimizes only execution time; wages are held fixed. We had this earlier too. Nothing about it is new.
  \item \emph{Medium run [new]:} the firm optimizes both skill and time costs. This is similar in flavor to the old job design problem, but without hand-offs or multiple workers. The medium run captures the idea that AI affects training and workers' wages in addition to reducing execution time.
  \end{itemize}
\item \textbf{Minimal pointers to the long run and the job design problem.} When introducing workers' wages and the firm's optimization, I kept brief references to the job design appendix, noting that in the long run job boundaries are also allowed to adjust endogenously and we do consider it, but not in the main body.
\end{itemize}

\section*{Discussion}
\begin{itemize}
\item Revised to add formalism to the implications, so the section becomes more than just illustrative examples. Previously, only the fragmentation-index portion provided a formal discussion, and the rest were just examples. 
\item \textbf{Fragmentation-index argument:} kept as is, with no changes.
\item \textbf{Removed} the earlier reskilling and upskilling discussion, which was tied to the job design margin.
\item \textbf{New subsection on the overturning of comparative advantage.} Previously, this point was made only informally; I added a couple of propositions, motivated by examples, to show the economic forces that drive the breakdown of comparative advantage under chaining. [I'm open to feedback and suggestions here. If you think I've gone overboard, open and happy to cut pieces down.]
\item \textbf{Non-monotone returns to AI improvement:} reworked so the point is driven purely by chaining within a single job rather than by job design and job boundaries (which is how we argued it previously). I brought a proposition showing that whenever a reorganization threshold (triggered by chaining alone, not job redesign) is crossed, the marginal returns to improving AI jump upward discretely. Also, earlier we informally claimed that as AI quality $\alpha$ rises, the marginal returns to AI improvement would become steeper, since we would have longer chains and thus larger powers of $1/\alpha$ in the cost function. I realized this claim is not quite accurate: we won't necessarily have longer chains as $\alpha$ rises, so I removed it.
\end{itemize}

\section*{Empirics}
\begin{itemize}
\item Essentially unchanged. I retained the EC version, which itself was a condensed summary of the NBER working paper with only the core pieces. All robustness tests and the dataset construction details are moved to the appendix.
\item \textbf{New result.} I also have a new result, which I haven't fully ported into the draft. A caveat of our empirics was that we used \textbf{all} the tasks that O*NET assigns to an occupation, even those performed with very low frequency (e.g., once a year). I ran our analyses restricting each occupation to only its frequently performed tasks (those carried out roughly daily or more often) and asked whether the findings survive on this narrower subsample. The answer is that they do. I still haven't finalized the write-up or how best to present these results, but just wanted to flag that they exist and we can add them to the appendix if needed.
\end{itemize}

\end{document}
